\documentclass{article}
\usepackage[backend=biber,natbib=true,style=authoryear]{biblatex}
\addbibresource{/home/nguyen/1_NQBH/math/bib.bib}
\usepackage[utf8]{inputenc}
\usepackage{tocloft}
\renewcommand{\cftsecleader}{\cftdotfill{\cftdotsep}} % for sections, if you really want! (It is default in report and book class (So you may not need it).
\usepackage{float}
\usepackage{graphicx}
\usepackage[colorlinks=true,linkcolor=blue,urlcolor=red,citecolor=magenta]{hyperref}
\usepackage{amsmath,amssymb,amsthm}
\allowdisplaybreaks
\numberwithin{equation}{section}
\newtheorem{assumption}{Assumption}
\newtheorem{lemma}{Lemma}
\newtheorem{corollary}{Corollary}
\newtheorem{definition}{Definition}
\newtheorem{proposition}{Proposition}
\newtheorem{theorem}{Theorem}
\newtheorem{notation}{Notation}
\newtheorem{remark}{Remark}
\newtheorem{example}{Example}
\newtheorem{ques}{Question}
\newtheorem{problem}{Problem}
\newtheorem{conjecture}{Conjecture}
\usepackage[left=0.5in,right=0.5in,top=1.5cm,bottom=1.5cm]{geometry}
\def\labelitemii{$\circ$}

\title{Don Miguel Ruiz. \textit{The Four Agreements: A Practical Guide to Personal Freedom}. A Toltec Wisdom Book}
\author{Collector: Nguyen Quan Ba Hong}
\date{\today}

\begin{document}
\maketitle
\setcounter{secnumdepth}{6}
\setcounter{tocdepth}{6}
\tableofcontents
\vspace{5mm}
\noindent
\textbf{Note.} The term ``black magic'' is not meant to convey racial connotation; it is merely used to describe the use of magic for adverse or harmful purposes.

\begin{quotation}
    To the \textit{Circle of Fire}; those who have gone before, those who are present, and those who have yet to come.
\end{quotation}

%------------------------------------------------------------------------------%

\section*{Acknowledgments}
\textsc{I would like to humbly acknowledge my} mother Sarita, who taught me unconditional love; my father Jose Luis, who taught me discipline; my grandfather Leonardo Macias, who gave me the key to unlock the Toltec mysteries; and my sons Miguel, Jose Luis, and Leonardo.

%
I wish to express my deep affection and appreciation to the dedication of Gaya Jenkins and Trey Jenkins.

%
I would like to extend my profound gratitude to Janet Mills - publisher, editor, believe.

I am also abidingly grateful to Ray Chambers for lighting the way.

%
I would like to honor my dear friend Gini Gentry, an amazing ``brain'' whose faith touched my heart.

%
I would like to pay tribute to the many people who have given freely of their time, hearts, and resources to support the teachings.

A partial list includes: $\ldots$ (see \textbf{[Ruiz2018]}).

\section*{The Toltec}
\textsc{Thousands of years ago, the Toltec were} known throughout southern Mexico as ``women and men of knowledge.''

Anthropologists have spoken of the Toltec as a nation or a race, but, in fact, the Toltec were scientists and artists who formed a society to explore and conserve the spiritual knowledge and practices of the ancient ones.

They came together as masters (\textit{naguals}) and students oat Teotihuacan, the ancient city of pyramids outside Mexico City known as the place where ``\textit{Man Becomes God}.''

%
Over the millennia, the \textit{naguals} were forced to conceal the ancestral wisdom and maintain its existence in obscurity.

European conquest, coupled with rampant misuse of personal power by a few of the apprentices, made it necessary to shield the knowledge from those who were not prepare to use it wisely or who might intentionally misuse it for personal gain.

%
Fortunately, the esoteric Toltec knowledge was embodied and passed on through generations by different lineages of \textit{naguals}.

Though it remained veiled in secrecy for hundreds of years, ancient prophecies foretold the coming of an age when it would be necessary to return the wisdom to the people.

Now, don Miguel Ruiz, a \textit{nagual} from the Eagle Knight lineage, has been guided to share with us the powerful teachings of the Toltec.

%
Toltec knowledge arises from the same essential unity of truth as all the sacred esoteric traditions found around the world.

Though it is not a religion, it honors all the spiritual masters who have taught on the earth.

While it does embrace spirit, it is most accurately described as a way of life, distinguished by the ready accessibility of happiness and love.

\section*{Introduction: The Smokey Mirror}
\textsc{3000 years ago, there was a human} just like you and me who lived near a city surrounded by mountains.

The human was studying to become a medicine man, to learn the knowledge of his ancestors, but he didn't completely agree with everything he was learning.

In his heart, he felt there must be something more.

%
1 day, as he slept in a cave, he dreamed that he saw his own body sleeping.

He came out of the cave on the night of a new moon.

The sky was clear, and he could see millions of stars.

Then something happened inside of him that transformed his life forever.

He looked at his hands, he felt his body, and he heard his own voice say, ``I am made of light; I am made of stars.''

%
He looked at the stars again, and he realized that it's not the stars that create light, but rather light that creates the stars.

``Everything is made of light,'' he said, ``and the space in-between isn't empty.''

And he knew that everything that exists in 1 \textit{living being}, and that light is the messenger of life, because it is alive and contains all information.

%
Then he realized that although he was made of stars, he was not those stars.

``I am in-between the stars,'' he thought.

So he called the stars the \textit{tonal} and the light between the stars the \textit{nagual}, and he knew that what created the harmony and space between the two is Life or Intent.

Without Life, the \textit{tonal} and the \textit{nagual} could not exist.

Life is the force of the absolute, the supreme, the Creator who creates everything.

%
This is what he discovered: Everything in existence is a manifestation of the one living being we call God.

\textit{Everything is God}.

And he came to the conclusion that human perception is merely light perceiving light.

He also saw that matter is a mirror - everything is a mirror that reflects light and creates images of that light - and the world of illusion, the \textit{Dream}, is just like smoke which doesn't allow us to see what we really are.

``The real us is pure love, pure light,'' he said.

%
This realization changed his life.

Once he knew what he really was, he looked around at other humans and the rest of nature, and he was amazed at what he saw.

He saw himself in everything - in every human, in every animal, in every tree, in the water, in the rain, in the clouds, in the earth.

And he saw that Life mixed the \textit{tonal} and the \textit{nagual} in different ways to create billions of manifestations of Life.

%
In those few moments he comprehended everything.

He was very excited, and his heart was filled with peace.

He could hardly wait to tell his people what he had discovered.

But there were no words to explain it.

He tried to tell the others, but they could not understand.

They would see that he had changed, that something beautiful was radiating from his eyes and his voice.

They noticed that he no longer had judgment about anything or anyone.

He was no longer like anyone else.

%
He could understand everyone very well, but no one could understand him.

They believed that he was an incarnation of God, and he smiled when he heard this and he said,
\begin{quotation}
    ``It is true.
    
    I am God.
    
    But you are also God.
    
    We are the same, you and I.
    
    We are images of light.
    
    We are God.''
\end{quotation}
But still the people didn't understand him.

%
He had discovered that he was a mirror for the rest of the people, a mirror in which he could see himself.

``Everyone is a mirror,'' he said.

He saw himself in everyone, but nobody saw him as themselves.

And he realized that everyone was dreaming, but without awareness, without knowing what they really are.

They couldn't see him as themselves because there was a wall of fog or smoke between the mirrors.

And that wall of fog was made by the interpretation of images of light - the \fbox{\textit{Dream} of humans}.

%
Then he knew that he would soon forget all that he had learned.

He wanted to remember all the visions he had had, so he decided to call himself the Smokey Mirror so that he would always know that matter is a mirror and the smoke in-between is what keeps us from knowing what we are.

He said, ``I am the Smokey Mirror, because I am looking at myself in all of you, but we don't recognize each other because of the smoke in-between us.

That smoke is the \textit{Dream}, and the mirror is you, the dreamer.''
\begin{quotation}
    ``\textit{Living is easy with eyes closed, misunderstanding all you see}$\ldots$'' - John Lennon
\end{quotation}

\section{Domestication \& the Dream of the Planet}
\textsc{What you are seeing and hearing right now is} nothing but a dream.

You are dreaming right now in this moment.

You are dreaming with the brain awake.

%
Dreaming is the main function of the mind, and the mind dreams 24 hours a day.

It dreams when the brain is awake, and it also dreams when the brain is asleep.

The difference is that when the brain is awake, there is a material frame that makes us perceive things in a linear way.

When we go to sleep we do not have the frame, and the dream has the tendency to change constantly.

%
Humans are dreaming all the time.

Before we were born the humans before us created a big outside dream that we will call society's dream or \textit{the dream of the planet}.

The dream of the planet is the collective dream of billions of smaller, personal dreams, which together create a dream of a family, a dream of a community, a dream of a city, a dream of a country, and finally a dream of the whole humanity.

The dream of the planet includes all of society's rules, its beliefs, its laws, its religions, its different cultures and ways to be, its governments, schools, social events, and holidays.

%
We are born with the capacity to learn how to dream, and the humans who live before us teach us how to dream the way society dreams.

The outside dream has so many rules that when a new human is born, we hook the child's attention and introduce these rules into his or her mind.

The outside dream uses Mom and Dad, the schools, the religion to teach us how to dream.

%
\fbox{\textit{Attention} is the ability we have to discriminate} and to focus only on that which we want to perceive.

We can perceive millions of things simultaneously, but using our attention, we can hold whatever we want to perceive in the foreground of our mind.

The adults around us hooked our attention and put information into our minds through repetition.

That is the way we learned everything we know.

%
By using our attention we learned a whole reality, a whole dream.

We learned how to behave in society: what to believe and what not to believe; what is acceptable and what is not acceptable; what is good and what is bad; what is beautiful and what is ugly; what is right and what is wrong.

It was all there already - all that knowledge, all those rules and concepts about how to behave in the world.

%
When you were in school, you sat in a little chair and put your attention on what the teacher was teaching you.

When you went to church, you put your attention on what the priest or minister was telling you.

It is the same dynamic with Mom and Dad, brothers and sisters: They were all trying to hook your attention.

We also learn to hook the attention of other humans, and we develop a need for attention which can become very competitive.

Children compete for the attention of their parents, their teachers, their friends.

``Look at me!

Look at what I'm doing!

Hey, I'm here.''

The need for attention becomes very strong and continues into adulthood.

%
The outside dream hooks our attention and teaches us what to believe, beginning with the language that we speak.

Language is the code for understanding and communication between humans.

Every letter, every word in each language is an agreement.

We call this a page in a book; the word \textit{page} is an agreement that we understand.

\fbox{\textit{Once we understand the code, our attention is hooked and the energy is transferred from 1 person to another}.}

%
It was not your choice to speak English.

You didn't choose your religion or your moral values - they were already there before you were born.

We never had the opportunity to choose what to believe or what not to believe.

We never chose even the smallest of these agreements.

We didn't even chose our own name.

%
As children, we didn't have the opportunity to choose our beliefs, but we \textit{agreed} with the information that was passed to us from the dream of the planet via other humans.

\fbox{\textit{The only way to store information is by agreement}.}

The outside dream may hook our attention, but if we don't agree, we don't store that information.

As soon as we agree, we \textit{believe} it, and this is called \textit{faith}.

\fbox{\textit{To have faith is to believe unconditionally}.}

%
That's how we learn as children.

Children believe everything adults say.

We agree with them, and our faith is so strong that \textit{the belief system controls our whole dream of life}.

We didn't choose these beliefs, and we may have rebelled against them, but we were not strong enough to win the rebellion.

The result is surrender to the beliefs with out \textit{agreement}.

%
I call this process \textit{the domestication of humans}.

And though this domestication we learn how to live and how to dream.

In human domestication, the information from the outside dream is conveyed to the inside dream, creating our whole belief system.

1st the child is taught the names of things: Mom, Dad, milk, bottle.

Day by day, at home, at school, at church, and from television, we are told how to live, what kind of behavior is acceptable.

\textit{The outside dream teaches us how to be a human}.

We have a whole concept of what a ``woman'' is and what a ``man'' is.

And we also learn to judge: We judge ourselves, judge other people, judge the neighbors.

%
Children are domesticated the same way that we domesticate a dog, a cat, or any other animal.

In order to teach a dog we punish the dog and we give it rewards.

We train our children whom we love so much the same way that we train any domesticated animal: with a system of punishment and reward.

We are told, ``You're a good boy,'' or ``You're a good girl,'' when we do what Mom and Dad want us to do.

When we don't, we are ``a bad girl'' or ``a bad boy.''

%
When we went against the rules we were punished; when we went along with the rules we got a reward.

We were punished many times a day, and we were also rewarded many times a day.

Soon we became afraid of being punished and also afraid of not receiving the reward.

The reward is the attention that we got from our parents or from other people like siblings, teachers, and friends.

\textit{We soon develop a need to hook other people's attention in order to get the reward}.

%
The rewards feels good, and we keep doing what others want us to do in order to get the reward.

With that fear of being punished and that fear of not getting the reward, we start pretending to be what we are not, just to please others, just to be good enough for someone else.

We try to please Mom and Dad, we try to please the teachers at school, we try to please the church, and so we start acting.

\textit{We pretend to be what we are not because we are afraid of being rejected}.

The fear of being rejected becomes the fear of not being good enough.

Eventually we become someone that we are not.

We become a copy of Mamma's beliefs, Daddy's beliefs, society's beliefs, and religion's beliefs.

%
All our normal tendencies are lost in the process of domestication.

And when we are old enough for our mind to understand, we learn the word \textit{no}.

The adults say, ``Don't do this and don't do that.''

We rebel and say, ``No!''

\fbox{\textit{We rebel because we are defending our freedom}.}

We want to be ourselves, but we are very little, and the adults are big and strong.

After a certain time we are afraid because we know that every time we do something wrong we are going to be punished.

%
The domestication is so strong that at a certain point in our lives we no longer need anyone to domesticate us.

We don't need Mom or Dad, the school or the church to domesticate us.

\textit{We are so well trained that we are our own domesticator}.

We are an autodomesticated animal.

We can now domesticate ourselves according to the same belief system we were given, and using the same system of punishment and reward.

We punish ourselves when we don't follow the rules according to our belief system; we reward ourselves when we are the ``good boy'' or ``good girl.''

%
The belief system is like a Book of Law that rules our mind.

Without question, whatever is in that Book of Law, is our truth.

We base all of our judgments according to the Book of Law, even if these judgments go against our own inner nature.

Even moral laws like the Ten Commandments are programmed into our mind in the process of domestication.

1 by 1, all these agreements go into the Book of Law, and these agreements rule our dream.

%
There is something in our minds that judges everybody and everything, including the weather, the dog, the cat - everything.

The \fbox{\textit{inner judges}} uses what is in our Book of Law to judge everything we do and don't do, everything we think and don't think, and everything we feel and don't feel.

Everything lives under the tyranny of this Judge.

Every time we do something that goes against the Book of Law, the Judge says we are guilty, we need to be punished, we should be ashamed.

This happens many times a day, day after day, for all the years of our lives.

%
There is another part of us that receives the judgments, and this part is called the Victim.

The Victim carries the blame, the guilt, and the shame.

It is the part of us that says, ``Poor me, I'm not good enough, I'm not intelligent enough, I'm not attractive enough, I'm not worthy of love, poor me.''

The bid Judge agrees and says, ``yes, you are not good enough.''

And this is all based on a belief system that we never chose to believe.

These beliefs are so strong, that even years later when we are exposed to new concepts and try to make our own decisions, we find that these beliefs still control our lives.

%
Whatever goes against the Book of Law will make you feel a funny sensation in your solar plexus, and it's called \textit{fear}.

Breaking the rules in the Book of Law opens your \textit{emotional wounds}, and \textit{your reaction is to create emotional poison}.

Because everything that is in the Book of Law has to be true, anything that challenges what you believe is going to make you \textit{feel unsafe}.

Even if the Book of Law is wrong, it makes you \textit{feel safe}.

%
That is why we need a great deal of courage to challenge our own beliefs.

Because even if we know we didn't choose all these beliefs, it is also true that we agreed to all of them.

The agreement is so strong that even if we understand the concept of it not being true, we feel the blame, the guilt, and the shame that occur if we go against these rules.

%
Just as the government has a book of laws that rule the society's dream, our belief system is the Book of Laws that rules our personal dream.

All these laws exist in our mind, we believe them, and the Judge inside us bases everything on these rules.

The Judge decrees, and the Victim suffers the guilt and punishment.

But who says there is justice in this dream?

True justice is paying only once for each mistake.

True \textit{injustice} is paying more than once for each mistake.

%
How many times do we pay for 1 mistake?

The answer is thousands of times.

The human is the only animal on earth that pays a thousand times for the same mistake.

The rest of the animals pay once for every mistake they make.

But not us.

We have a powerful memory.

We make a mistake, we judge ourselves, we find ourselves guilty, and we punish ourselves.

If justice exists, then that was enough; we don't need to do it again.

But every time we remember, we judge ourselves again, we are guilty again, and we punish ourselves again, and again, and again.

If we have a wife or husband he or she also reminds us of the mistake, so we can judge ourselves again, punish ourselves again, and find ourselves guilty again.

\textit{Is this fair?}

%
How many times do we make our spouse, our children, or our parents pay for the same mistake?

Every time we remember the mistake, we blame them again and send them all the emotional poison we feel at the justice, and then we make them pay again for the same mistake.

\textit{Is that justice?}

The Judge in the mind is wrong because the belief system, the Book of Law, is wrong.

The whole dream is based on false law.

95\% of the beliefs we have stored in our minds are nothing but lies, and we suffer because we believe al these lies.

%
In the dream of the planet it is normal for humans to suffer, to live in fear, and to create \textit{emotional dramas}.

The outside dream is not a pleasant dream; it is a dream of violence, a dream of fear, a dream of war, a dream of injustice.

The personal dream of humans will vary, but globally it is mostly a nightmare.

If we look at human society we see a place so difficult to live in because it is ruled by fear.

Throughout the world we see human suffering, anger, revenge, addictions, violence in the street, and tremendous injustice.

It may exist at different levels in different counties around the world, but fear is controlling the outside dream.

%
If we compare the dream of human society with the description of hell that religions all around the world have promulgated, we find they are exactly the same.

Religions say that hell is a place of punishment, a place of fear, pain, and suffering, a place where the fire burns you.

Fire is generated by emotions that come from fear.

Whenever we feel the emotions of anger, jealousy, envy, or hate, we experience a fire burning within us.

We are living in a dream of hell.

%
\textit{If you consider hell as a state of mind, the hell is all around us}.

Others may warn us that if we don't do what they say we should do, we will go to hell.

Bad news!

\fbox{\textit{We are already in hell, including the people who tell us that}.}

No human can condemn another to hell because we are already there.

Others can put us into a deeper hell, true.

But only if we allow this to happen.

%
Every human has his or her own personal dream, and just like the society dream, it is \textit{often ruled by fear}.

We learn to dream hell in our own life, in our personal dream.

The same fears manifest in different ways for each person, of course, but we experience anger, jealousy, hate, envy, and other negative emotions.

Our personal dream can also become an ongoing nightmare where we suffer and live in a state of fear.

But we don't need to dream a nightmare.

It is possible to enjoy a pleasant dream.

%
All of humanity is searching for truth, justice, and beauty.

We are on an eternal search for the truth because we only believe in the lies we have stored in our mind.

We are searching for justice because in the belief system we have, there is no justice.

We search for beauty because it doesn't matter how beautiful a person is, we don't believe that person has beauty.

We keep searching and searching, when everything is already within us.

There is no truth to find.

Wherever we turn our heads, all we see is the truth, but with the agreements and beliefs we have stored in our mind, we have no eyes for this truth.

%
\textit{We don't see the truth because we are blind}.

\textit{What blinds us are all those false beliefs we have in our mind}.

We have the need to be right and to make others wrong.

We trust what we believe, and our beliefs set us up for suffering.

It is as if we live in the middle of a fog that doesn't let us see any further than our own nose.

We live in a fog that is not even real.

This fog is a dream, your personal dream of life - what you believe, all the concepts you have about what you are, all the agreements you have made with others, with yourself, and even with God.

%
Your whole mind is a fog which the Toltecs called a \textit{mitote} (pronounced MIH-TOE' -TAY).

\fbox{\textit{Your mind is a dream where a thousand people at the same time, and nobody understands each other}.}

This is the condition of the human kind - a big \textit{minote}, and with that big \textit{mitote} you cannot see what you really are.

In India they call the \textit{mitote maya}, which means \textit{illusion}.

It is the personality's notion of ``I am.''

Everything you believe about yourself and the world, all the concepts and programming you have in your mind, are all the \textit{minote}.

\textit{We cannot see who we truly are; we cannot see that we are not free}.

%
\textit{That is why humans resist life}.

\textit{To be alive is the biggest fear humans have}.

Death is not the biggest fear we have; \textit{our biggest fear is taking the risk to be alive - the risk to be alive and express what we really are}.

Just being ourselves is the biggest fear of humans.

We have learned to live our lives trying to satisfy other people's demands.

We have learned to live by other people's points of view because of the fear of not being accepted and of not being good enough for someone else.

%
During the process of domestication, we form an image of what perfection is in order to try to be good enough.

We create an image of how we should be in order to be accepted by everybody.

We especially try to please the ones who love us, like Mom and Dad, big brothers and sisters, the priests and the teacher.

Trying to be good enough for them, we create an image of perfection, but we don't fit this image.

We create this image, but this image is not real.

We are never going to be perfect from this point of view.

\textit{Never!}

%
\fbox{\textit{Not being perfect, we reject ourselves}.}

And the level fo self-rejection depends upon how effective the adults were in breaking our integrity.

After domestication it is no longer about being good enough for anybody else.

We are not good enough for ourselves because we don't fit with our own image of perfection.

We cannot forgive ourselves for not being what we wish to be, or rather what we \textit{believe} we should be.

We cannot forgive ourselves for not being perfect.

%
We know we are not what we believe we are supposed to be and so we feel false, frustrated, and dishonest.

We try to hide ourselves, and we pretend to be what we are not.

The result is that we feel unauthentic and wear social masks to keep others from noticing this.

We are so afraid that somebody else will notice that we are not what we pretend to be.

We judge others according to our image of perfection as well, and naturally they fall short of our expectations.

%
\textit{We dishonor ourselves just to please other people}.

We even do harm to our physical bodies just to be accepted by others.

You see teenagers taking drugs just to avoid being rejected by other teenagers.

They are not aware that the problem is that they don't accept themselves.

They reject themselves because they are not what they pretend to be.

They wish to be a certain way, but they are not, and for this they carry shame and guilt.

Humans punish themselves endlessly for not being what they believe they should be.

They become very self-abusive, and they use other people to abuse themselves as well.

%
\textit{But nobody abuses us more than we abuse ourselves}, and it is the Judge, the Victim, and the belief system that make us do this.

True, we find people who say their husband or wife, or mother or father, abused them, but you know that we abuse ourselves much more than that.

\textit{The way we judge ourselves is the worst judge that ever existed}.

If we make a mistake in front of people, we try to deny the mistake and cover it up.

But as soon as we are alone, the Judge becomes so strong, the guilt is so strong, and we feel so stupid, or so bad, or so unworthy.

%
\textit{In your whole life nobody has ever abused you more than you have abused yourself}.

\fbox{\textit{And the limit of your self-abuse is exactly the limit that you will tolerate from someone else}.}

\textit{If someone abuses you a little more than you abuse yourself, you will probably walk away from that person}.

\textit{But if someone abuses you a little less than you abuse yourself, you will probably stay in the relationship and tolerate it endlessly}.

%
If you abuse yourself very badly, you can even tolerate someone who beats you up, humiliates you, and treats you like dirt.

\textit{Why?}

Because in your belief system you say, ``I deserve it.

This person is doing me a favor by being with me.

I'm not worthy of love and respect.

I'm not good enough.''

%
\textit{We have the need to be accepted and to be loved by others, but we cannot accept and love ourselves}.

The more self-love we have, the less we will experience self-abuse.

\textit{Self-abuse comes from self-rejection, and self-rejection comes from having an image of what it means to be perfect and never measuring up to that ideal}.

Our image of perfection is the reason we reject ourselves; it is why we don't accept ourselves the way we are, and why we don't accept others the way they are.
\begin{center}
    \textsc{Prelude to a New Dream}
\end{center}
There are thousands of agreements you have made with yourself, with other people, with your dream of life, with God, with society, with your parents, with your spouse, with your children.

But the most important agreements are the ones you made with yourself.

In these agreements you tell yourself \textit{who you are, what you feel, what you believe}, and \textit{how to behave}.

The result is what you call \textit{your personality}.

In these agreements you say, ``This is what I am.

This is what I believe.

I can do certain things, and some things I cannot do.

This is reality, that is fantasy; this is possible, that is impossible.''

%
1 single agreement is not such a problem, but we have many agreements that make us suffer, that make us fail in life.

If you want to live a life of joy and fulfillment, you have to find the courage to break those agreements that are fear-based and claim your \textit{personal power}.

The agreements that come from fear require us to expend a lot of energy, but the agreements that come from love help us to conserve energy and even gain extra energy.

%
Each of us is born with a certain amount of personal power that we rebuild every day after we rest.

Unfortunately, we spend all our personal power 1st to create all these agreements and then to keep these agreements.

Our personal power is dissipated by all the agreements we have created, and the result is that we feel powerless.

We have just enough power to survive each day, because most of it is used to keep the agreements that trap us in the dream of the planet.

\textit{How can we change the entire dream of our life when we have no power to change even the smallest agreement?}

%
If we can see it is our agreements that rule our own life, and we don't like the dream of our life, we need to change the agreements.

When we are finally ready to change our agreements, there are 4 very powerful agreements that will help us to break those agreements that come from fear and deplete our energy.

%
Each time you break an agreement, all the power you used to create it returns to you.

If you adopt these 4 new agreements, they will create enough personal power for you to change the entire system of your old agreements.

%
You need a very strong will in order to adopt The Four Agreements - but if you can begin to live your life with these agreements, the transformation in your life will be amazing.

You will see the drama of hell disappear right before your very eyes.

Instead of living in a dream of hell, you will be creating a new dream - your personal dream of heaven.

\section{\textsc{The 1st Agreement}: Be Impeccable with your Word}

\section{\textsc{The 2nd Agreement}: Don't Take Anything Personally}

\section{\textsc{The 3rd Agreement}: Don't Make Assumptions}

\section{\textsc{The 4th Agreement}: Always Do Your Best}

\section{\textsc{The Toltec Path to Freedom}: Breaking Old Agreements}

\section{\textsc{The New Dream}: Heaven on Earth}

\section{Prayers}

%------------------------------------------------------------------------------%

\printbibliography[heading=bibintoc]
\end{document}